% --------------------
%  table of contents
% --------------------
\usepackage{tocloft}
\usepackage{setspace}

% Вертикальный воздух между пунктами TOC.
% Внутри пункта строки идут одинарным интервалом (см. \fontsize ниже),
% а между пунктами добавляется небольшой пробел, имитирующий полуторный.
\setlength{\cftbeforesecskip}{0pt}
\setlength{\cftbeforesubsecskip}{0pt}
\setlength{\cftbeforesubsubsecskip}{0pt}

% отступы:
% > разделы 0
% > подразделы 0.49см
% > подпункты 0.99см
\setlength{\cftsecindent}{0pt}
\setlength{\cftsubsecindent}{0.49cm}
\setlength{\cftsubsubsecindent}{0.99cm}

% Ширина колонки номера
\setlength{\cftsecnumwidth}{0.75em}
\setlength{\cftsubsecnumwidth}{1.55em}
\setlength{\cftsubsubsecnumwidth}{2.7em}

% Плотность отточия
\renewcommand{\cftdotsep}{1.1}

% Ширина колонки номера страницы (2 цифры в кегле 14pt помещаются в 1.3em)
\cftsetpnumwidth{0.85em}
\cftsetrmarg{0.82em}

\renewcommand{\cfttoctitlefont}{\hfill\fontsize{14pt}{\baselineskip}\selectfont\bfseries}
\renewcommand{\cftaftertoctitle}{\hfill}
\renewcommand{\cftaftertoctitleskip}{1em}

% Отточие идёт вплотную к номеру страницы (без принудительного \hspace)
\renewcommand{\cftsecleader}{\cftdotfill{\cftdotsep}}
\renewcommand{\cftsubsecleader}{\cftdotfill{\cftdotsep}}
\renewcommand{\cftsubsubsecleader}{\cftdotfill{\cftdotsep}}

\addto\captionsrussian{%
	\renewcommand{\contentsname}{СОДЕРЖАНИЕ}%
}

% Шрифт записей TOC: одинарный интервал (baselineskip = 17pt при кегле 14pt).
% Так многострочные пункты не разъезжаются.
\renewcommand{\cftsecfont}{\normalfont\fontsize{14pt}{14pt}\selectfont}
\renewcommand{\cftsubsecfont}{\normalfont\fontsize{14pt}{14pt}\selectfont}
\renewcommand{\cftsubsubsecfont}{\normalfont\fontsize{14pt}{14pt}\selectfont}

% Шрифт номеров страниц: тот же одинарный интервал, без принудительной ширины
\renewcommand{\cftsecpagefont}{\normalfont\fontsize{14pt}{14pt}\selectfont}
\renewcommand{\cftsubsecpagefont}{\normalfont\fontsize{14pt}{14pt}\selectfont}
\renewcommand{\cftsubsubsecpagefont}{\normalfont\fontsize{14pt}{14pt}\selectfont}

% Гасим растягивание пробелов в многострочных пунктах TOC.
% tocloft переопределяет \l@section/\l@subsection/\l@subsubsection своими
% макросами, которые ставят \rightskip \@tocrmarg (фиксированное правое поле)
% и тем самым justify-ят все строки заголовка. Через etoolbox\patchcmd
% заменяем правое поле на упругое \@tocrmarg plus 1fil — теперь нерастянутые
% строки оставляют пустое пространство справа, без растягивания пробелов.
\usepackage{etoolbox}
\makeatletter
\patchcmd{\l@section}
  {\rightskip \@tocrmarg}{\rightskip \@tocrmarg plus 1fil}{}{}
\patchcmd{\l@subsection}
  {\rightskip \@tocrmarg}{\rightskip \@tocrmarg plus 1fil}{}{}
\patchcmd{\l@subsubsection}
  {\rightskip \@tocrmarg}{\rightskip \@tocrmarg plus 1fil}{}{}
\makeatother
